\chapter{AGI 培育工程：课程、闭环与长期能力形成}
\label{chap:agi_cultivation}

本章讨论的不是一次训练作业怎样把损失压到最低，而是开放智能系统怎样在长期运行中形成能力、校正边界并承担行动后果。我把这种持续过程称为“培育”，因为模型参数、外部记忆、工具权限、社会关系与评估制度会共同改变系统的后续轨迹；任何一次更新都只是其中的局部事件。HSF-HD在这里提供一般状态动力学：它描述状态怎样受任务、输入、结构与环境反馈共同推动；AgentOS则是可替换的工程实例，用 $h(t)/E/\Theta$、$\Psi$、SPC、TCE、CCM、Boundary、Offloading 与外部记忆 $M^e$ 承载观测、控制和审计。二者不能混同，更不能由某个实现有效便宣布所有智能都必须采用相同部件。原章使用“相变、材料、钻孔”等物理语言，本章保留这些词作为直观索引，却把真正的论证落在可定义的数据类型、干预实验、阶段门槛和故障恢复上。所谓成熟也不等于某一标量达到峰值，而是系统在能力、诚实、边界、协作和可逆性之间形成可持续的联合状态。

\section{培育对象、数据次序与证据边界}
\label{sec:adiabatic_paving_and_nert_spectrometer}

先定义培育对象。离散检查时刻记为 $k\in\mathbb N$，间隔 $\Delta t_k>0$ 可以是一次训练步、一次部署回合或一次人工复盘周期。系统状态写成
\begin{equation}
R_k=(h_k,E_k,\Theta_k,\Psi_k,M^e_k,q_k,\nu_k)\in
\mathcal H\times\mathcal E\times\mathcal T\times\mathcal P\times\mathcal M\times\mathcal Q\times\mathcal V .
\end{equation}
$h_k$ 是内部表征的受控读出，$E_k$ 是可用算力、时间和调用预算，$\Theta_k$ 是目标、策略及模型版本，$\Psi_k$ 是正在执行的任务上下文，$M^e_k$ 是可持久化外部记忆，$q_k$ 是能力与风险评估向量，$\nu_k$ 是数据、工具和制度环境的版本。除实际资源项外，这些量通常是无量纲评分或类型化对象，不能与焦耳能耗混算。训练样本也不是裸文本，而是带来源与用途的记录 $z=(x,y,c,s,\tau,\rho)$：输入 $x$、监督或反馈 $y$、任务条件 $c$、来源 $s$、时间戳 $\tau$、许可和风险标签 $\rho$。培育控制 $u_k$ 包括采样分布、学习率、工具权限、记忆写入规则及人工复核强度，于是一般更新只写为
\begin{equation}
R_{k+1}=F_{\nu_k}(R_k,u_k,z_k,\xi_k),
\end{equation}
其中 $\xi_k$ 表示没有被模型显式控制的随机性。这个式子是建模接口，不是关于心智本体的定理。

原章把混合批次直接解释为“剪切应力”和幻觉的物理根源，这个结论过强。随机混合有时能降低梯度相关性，也可能在前置技能不足、任务冲突或标签制度不一致时造成负迁移。可检验的说法应当是：若两个数据簇 $D_a,D_b$ 在相同检查点上产生的归一化更新方向 $g_a,g_b$ 满足 $\langle g_a,g_b\rangle/(\|g_a\|\|g_b\|)<-\gamma$，并且交替训练使保留集性能持续下降，那么我们获得了任务干扰的证据；若只看到一次梯度尖峰，则还不能区分脏数据、数值不稳定和真正的概念冲突。课程学习因此不是“永远先轻后重”的教条，而是一个带反事实比较的调度问题。基础语法通常应早于长链证明，但故事中提前出现生死、道德或宇宙等复杂主题，也可能以低要求方式建立未来检索锚点。原章“童话掺杂重概念”的例子在这个意义上仍然成立：它说明难度至少有概念跨度、推理深度、表达密度和行动风险四个轴，而不是一条神秘的质量轴。

我把 Nert 保留为工程代号，定义为相对于当前状态的需求向量 $n(z\mid R_k)=(d,p,r,a)$。$d$ 是依赖深度，$p$ 是单位上下文的推理密度，$r$ 是错误后果等级，$a$ 是与既有能力的失配度；各分量由独立量表归一到 $[0,1]$，不可直接解释为物理质量。调度器根据验证过的先修图、历史迁移矩阵和风险预算选择样本，而非仅按 $d$ 排序。证据门槛也要写清：预测某批数据会带来收益只是模型假设；随机种子复现、消融对照和延迟保留测试才构成经验支持；由此得到的只是特定模型和版本上的课程效果，不足以推出内部结构与人脑、晶体或流形整体同构。

调度实验要把顺序效应与样本构成效应分开。三条训练轨迹使用相同的数据多重集和总计算预算：轨迹甲按先修图排序，轨迹乙随机混合，轨迹丙按逆序排列；每条轨迹在相同检查点测量即时任务、先修保留、跨域迁移和校准。只有当甲在多个种子上持续优于乙和丙，且差异不能由批尺寸、长度或优化器状态解释时，我们才说课程次序在该范围内有效。若甲只更快降低训练损失，却没有更好的延迟保留，则“绝热铺设”只是优化加速，不是能力结构改善。更换模型规模、分词器或检索组件后，结论必须重新验证。培育对象还包括制度状态：同一模型在只读问答、可执行代码和可支配资金的环境中具有完全不同的风险，故 $\nu_k$ 必须保存权限矩阵、审批链和责任归属。报告分别描述模型可生成什么、运行时允许什么、组织实际上部署什么，避免把潜在能力与已授权行动混淆。

\section{数据材料表征与任务相关阻抗}
\label{sec:cognitive_materials_science_nert_density}

数据确实具有不同训练作用，但“材料”在本章只是把多维响应组织起来的计算模型。对数据单元 $B$，我定义描述向量
\begin{equation}
m(B;R_k)=(\ell_B,d_B,p_B,n_B,c_B,r_B,b_B)\in[0,1]^7,
\end{equation}
依次表示许可与来源可信度、依赖深度、推理密度、新颖度、与当前任务的覆盖度、行动风险和可能的偏置强度。每一维都必须给出测量规程。例如推理密度可由完成任务所需的最短可验证步骤数除以有效上下文长度，再按语料域内分位数归一；新颖度可由当前模型对内容的预测残差与最近邻距离联合估计；风险不是文本“沉重”的修辞，而是错误输出经工具或用户链路传播后的最大可恢复损失等级。向量的单位是量表分值，不能与 token 数、梯度范数或电力消耗相加。原章所谓总质量、压力和峰值概念，分别可以映射为覆盖总量、局部要求密度和单点最大难度，但这些统计量不足以唯一决定训练后果。

相对于状态 $R_k$ 的“阻抗”可用预测模型表达。令 $\delta q_{k:H}(B)$ 表示在固定预算下引入 $B$ 后，未来 $H$ 个检查点的能力—风险变化向量；令 $C_k(B)$ 是训练、标注、复核和可能回滚的资源成本。我们用历史干预数据拟合
\begin{equation}
\widehat{\delta q}_{k:H}=G_\omega(m(B;R_k),q_k,\nu_k),\qquad
Z_k(B)=\Pr(\delta q^{\rm safe}_{k:H}<0)+\lambda_C\widetilde C_k(B),
\end{equation}
其中第一项是安全保留指标下降的校准概率，第二项是归一化成本。$Z_k$ 只是风险—成本指数，不是自然界常数。模型版本、评测集或工具边界变化时，$G_\omega$ 也必须重新校准。若高难数据与先修缺口同时出现，调度器可以拆分证明、补充例子、延长上下文或先训练中间任务；若数据表达简单却带高行动风险，则应提高人工复核而非误判为“轻数据”。这样才能保留原章“密度比总量更重要”的合理洞见，同时避免把长文本天然当作缓冲海绵。

一个反例说明为什么不能机械地先易后难。模型学习药物剂量换算时，通俗病例叙述很长、语言简单，却可能隐藏单位换算错误；短小的量纲分析题反而更容易验证。按上下文长度稀释“认知压力”会把前者排到安全区，实际却提高风险。正确做法是把单位、药物、患者条件和允许工具写入类型约束，要求计算结果通过独立计算器和范围检查，再比较不同课程对保留集的影响。另一个反例是代码库：致密源码不一定比教程更难，如果模型已有相应 API 结构；难度必须相对于 $R_k$ 而言。数据相图不应画成固定的物理相区，而应表示任务相关的收益、干扰、成本和风险前沿。

课程选择可以写成有限时域优化：在候选集合 $\mathcal B_k$ 中选择分布 $\pi_k$，最大化经不确定性惩罚的预期能力增量，同时满足许可、风险、遗忘和预算约束。若各指标已归一，可用加权目标作为工程近似；权重必须记录，且目标下降不等于真实任务成功。每个周期保留跨阶段探针，检查系统是否只适应了调度器自身的分类。所谓“拓扑播种”于是被重述为低行动风险的早期概念暴露；“凝结”是补齐可验证依赖链；“造山”是对长链和跨域任务进行受控压力测试。它们是课程设计策略，不是必然发生的热力学相变。

材料的组合也不是各分量线性相加。两批单独安全的数据合并后，可能因指令优先级冲突、实体同名或时序倒置产生新错误；两批单独困难的数据若分别提供问题和验证器，组合后反而更容易学习。调度器应抽样测量二阶交互，至少比较 $B_i$、$B_j$ 与 $B_i\cup B_j$ 的影响轨迹，并把差异记入组合风险。长文档还要区分内容密度与检索负担：模型未找到关键段落，不代表不理解段落本身。可以用带位置提示、无提示和错误提示三组测试分离检索、绑定与推理失败。多模态材料则检查图像、文字和元数据是否指向同一对象，避免模型依靠单一通道捷径。若材料评分不能预测真实响应，应废弃评分体系，而不是把失败归结为“材料尚未成熟”。

\section{测量系统、影响轨迹与因果校准}
\label{sec:cognitive_mass_spectrometer_engineering}

认知质谱仪可以留下这个名字，但它必须被实现成测量流水线而不是权威分类器。流水线包含版本化探针模型、受控干预器、材料响应编码器和在线校准器。探针模型接收同一数据单元的原版、删减版、顺序置换版与事实反转版；干预器保证除被研究因素外，其余训练预算、初始检查点和随机种子尽量一致；编码器把响应压缩为可比较统计量；校准器用目标模型的真实后续表现修正代理偏差。所有输出绑定模型哈希、数据哈希、优化器版本和评测版本。没有这些版本信息，同一分数不能跨实验直接比较。

影响轨迹不定义为不可计算的 $\partial\theta/\partial x$，而定义为有限干预差。对检查点参数 $\theta_k\in\mathbb R^P$，以相同步长训练 $J$ 步，分别得到包含数据单元 $B$ 的 $\theta^{+}_{k,J}$ 和以匹配对照替换 $B$ 的 $\theta^{-}_{k,J}$。在读出函数 $f_j:\mathbb R^P\to\mathbb R$ 上记录
\begin{equation}
I_j(B;k,J)=f_j(\theta^{+}_{k,J})-f_j(\theta^{-}_{k,J}),\quad j=1,\ldots,m .
\end{equation}
$f_j$ 可以是能力准确率、校准误差、拒答质量、记忆泄漏率或梯度稳定性。实验在多个配对随机种子上重复并报告置信区间。梯度范数、Hessian 向量积和 Fisher 对角近似可作为诊断特征，但它们只是当前损失几何的局部量；梯度大不自动等于新知识，Fisher 大也不等于“曲率撕裂”。GPU 电耗、优化代价、信息熵和活动强度更不能称作同一种能量。

测量还必须区分可逆适应与持久结构改变。若 $B$ 只让训练损失快速下降，而在延迟保留、改写提问和分布外任务上没有收益，我把它标为短时拟合；若收益在撤除同类样本后仍保留，并能迁移到依赖任务，才标为持久学习。若能力提升伴随其他域显著下降，则记录迁移矩阵 $T_{ab}$，其中 $T_{ab}$ 表示在域 $a$ 的干预对域 $b$ 保留指标的标准化影响。原章的“做功与热耗散”可以作为直觉说明这一区分，但工程上使用的是持久收益和无效更新，不宣称它们对应物理功与热。影响轨迹也不能证明模型内部已经形成某种唯一概念结构，只能说明干预与选定读出之间的差异。

校准器面对三个常见陷阱。第一，小探针模型的响应排序未必能转移到目标模型，需在目标检查点抽样复测并计算秩相关与校准曲线。第二，评测污染会把记忆再现误当能力形成，保留集必须按来源隔离并加入新生成的结构等价题。第三，调度器会改变数据分布，使历史模型在新区域外推失真；输入漂移超过阈值时应回退到保守采样和人工审查。报告同时给出效应量、不确定区间和失败样本。若某批形式证明数据提高数学保留，却令自然语言解释的校准变差，就展示交换发生的任务与条件，而不能压成单一“高质量”标签。若改善来自更多计算而非材料本身，用等计算量对照抵消；若无法控制条件，则明确标记相关性经验主张。

测量流水线本身也会影响被测对象。频繁用同一探针选择训练数据，会使模型针对探针优化，原本独立的测量逐渐成为训练目标。团队应保存冻结的审计探针，并用轮换探针和人工新题检测适应；探针结果进入调度器时记录用途，避免把验证集悄然变成训练集。对无法公开的安全评测，可以公开任务类型、评分规则和治理过程，而保留具体样本。可重复性不要求泄露所有敏感材料，但要求独立评审者能够确认版本、控制变量和结论范围。质谱仪最重要的产品不是漂亮分数，而是一条从数据版本、干预设计、状态变化到复现实验的证据链。

% DRAFT BREAKPOINT: next section is sec:genesis_phase; sections 4--10 not yet reconstructed.
